% GENERATED FILE. Edit canonical lessons, then run npm run generate:books.
% canonical-source-hash: fnv1a64:82ffa0d02caf3d7f
% canonical-lessons: FR-C216-autour, FR-C216-partout, FR-C216-ailleurs, FR-C216-la-haut, FR-C216-banlieue

\chapter{Around, Everywhere, Elsewhere, Up there, Suburbs}
\label{ch:fr-around-everywhere-elsewhere-up-there-suburbs}
\hlchaptermodality{216}

\begin{chapteropening}
\textbf{By the end of this chapter:} \emph{I can say around, everywhere, elsewhere, up there and the suburbs.}

Complete the last lesson of chapter 216: I can say around, everywhere, elsewhere, up there and the suburbs.
\end{chapteropening}

\section[autour]{autour — around}

\label{lesson:FR-C216-autour}



\begin{quote}
Before the new one: say the French for a syrup, then the French for an ointment.
\end{quote}

\subsection*{You'll want to know: autour}
\textbf{autour} — \textquotedblleft{}around\textquotedblright{}.

\begin{cousinweb}
\textbf{autour} — \emph{au} and the word for \textquotedblleft{}a turn\textquotedblright{}, $\leftarrow$ Latin \emph{tornus}, \textquotedblleft{}a lathe\textquotedblright{}: round about. (EtymDB 2.1)
\end{cousinweb}

\begin{grammarlens}[title={What you've built}]
One more word for this chapter.
\end{grammarlens}

\subsection*{Your turn}
\begin{itemize}
\raggedright
  \item \emph{Say it:} \emph{autour}
  \item \emph{Say it:} \emph{autour}, once more
  \item \emph{Say it:} say \emph{autour}
  \item \emph{Recall:} say the French for a syrup, then the French for an ointment, then say \emph{autour} again
\end{itemize}

\begin{tcolorbox}[breakable,colback=teal!4,colframe=teal!35!black,title={Before you move on}]
What does autour mean? (\textquotedblleft{}Around\textquotedblright{}.) Say it once more.
\end{tcolorbox}
\section[partout]{partout — everywhere}

\label{lesson:FR-C216-partout}



\begin{quote}
Before the new one: say the French for an ointment, then the French for around.
\end{quote}

\subsection*{You'll want to know: partout}
\textbf{partout} — \textquotedblleft{}everywhere\textquotedblright{}.

\begin{cousinweb}
\textbf{partout} — the words for \textquotedblleft{}through\textquotedblright{} and \textquotedblleft{}all\textquotedblright{}, the second $\leftarrow$ Latin \emph{totus} (EtymDB 2.1): through all places.
\end{cousinweb}

\begin{grammarlens}[title={What you've built}]
One more word for this chapter.
\end{grammarlens}

\subsection*{Your turn}
\begin{itemize}
\raggedright
  \item \emph{Say it:} \emph{partout}
  \item \emph{Say it:} \emph{partout}, once more
  \item \emph{Say it:} say \emph{partout}
  \item \emph{Recall:} say the French for an ointment, then the French for around, then say \emph{partout} again
\end{itemize}

\begin{tcolorbox}[breakable,colback=teal!4,colframe=teal!35!black,title={Before you move on}]
What does partout mean? (\textquotedblleft{}Everywhere\textquotedblright{}.) Say it once more.
\end{tcolorbox}
\section[ailleurs]{ailleurs — elsewhere}

\label{lesson:FR-C216-ailleurs}



\begin{quote}
Before the new one: say the French for around, then the French for everywhere.
\end{quote}

\subsection*{You'll want to know: ailleurs}
\textbf{ailleurs} — \textquotedblleft{}elsewhere\textquotedblright{}.

\begin{cousinweb}
\textbf{ailleurs} $\leftarrow$ Old French \emph{aillors} (EtymDB 2.1) $\leftarrow$ Latin \emph{aliorsum}, \textquotedblleft{}toward another place\textquotedblright{} (Lewis \& Short) — inherited.
\end{cousinweb}

\begin{grammarlens}[title={What you've built}]
One more word for this chapter.
\end{grammarlens}

\subsection*{Your turn}
\begin{itemize}
\raggedright
  \item \emph{Say it:} \emph{ailleurs}
  \item \emph{Say it:} \emph{ailleurs}, once more
  \item \emph{Say it:} say \emph{ailleurs}
  \item \emph{Recall:} say the French for around, then the French for everywhere, then say \emph{ailleurs} again
\end{itemize}

\begin{tcolorbox}[breakable,colback=teal!4,colframe=teal!35!black,title={Before you move on}]
What does ailleurs mean? (\textquotedblleft{}Elsewhere\textquotedblright{}.) Say it once more.
\end{tcolorbox}
\section[là-haut]{là-haut — up there}

\label{lesson:FR-C216-la-haut}



\begin{quote}
Before the new one: say the French for everywhere, then the French for elsewhere.
\end{quote}

\subsection*{You'll want to know: là-haut}
\textbf{là-haut} — \textquotedblleft{}up there\textquotedblright{}.

\begin{cousinweb}
\textbf{là-haut} — \emph{là}, \textquotedblleft{}there\textquotedblright{}, and \emph{haut}, \textquotedblleft{}high\textquotedblright{}. \emph{Haut} $\leftarrow$ Latin \emph{altus}, \textquotedblleft{}high\textquotedblright{}, with an \emph{h} from Frankish (EtymDB 2.1 records the Frankish \emph{hauh}); \emph{là} $\leftarrow$ Latin \emph{illāc}, \textquotedblleft{}there\textquotedblright{}.
\end{cousinweb}

\begin{grammarlens}[title={What you've built}]
One more word for this chapter.
\end{grammarlens}

\subsection*{Your turn}
\begin{itemize}
\raggedright
  \item \emph{Say it:} \emph{là-haut}
  \item \emph{Say it:} \emph{là-haut}, once more
  \item \emph{Say it:} say \emph{là-haut}
  \item \emph{Recall:} say the French for everywhere, then the French for elsewhere, then say \emph{là-haut} again
\end{itemize}

\begin{tcolorbox}[breakable,colback=teal!4,colframe=teal!35!black,title={Before you move on}]
What does là-haut mean? (\textquotedblleft{}Up there\textquotedblright{}.) Say it once more.
\end{tcolorbox}
\section[la banlieue]{la banlieue — the suburbs}

\label{lesson:FR-C216-banlieue}



\begin{quote}
Before the new one: say the French for elsewhere, then the French for up there.
\end{quote}

\subsection*{You'll want to know: la banlieue}
\textbf{la banlieue} — \textquotedblleft{}the suburbs\textquotedblright{}.

\begin{cousinweb}
\textbf{banlieue} $\leftarrow$ Medieval Latin \emph{banleuca}: \emph{ban}, \textquotedblleft{}a lord's rule\textquotedblright{}, a Frankish word, and \emph{leuca}, \textquotedblleft{}a league\textquotedblright{}, a Gaulish one — the land within a league of a town. (EtymDB 2.1)
\end{cousinweb}

\begin{grammarlens}[title={What you've built}]
One more word for this chapter.
\end{grammarlens}

\subsection*{Your turn}
\begin{itemize}
\raggedright
  \item \emph{Say it:} \emph{la banlieue}
  \item \emph{Say it:} \emph{la banlieue}, once more
  \item \emph{Say it:} say \emph{la banlieue}
  \item \emph{Recall:} say the French for elsewhere, then the French for up there, then say \emph{la banlieue} again
\end{itemize}

\begin{tcolorbox}[breakable,colback=teal!4,colframe=teal!35!black,title={Before you move on}]
What does la banlieue mean? (\textquotedblleft{}The suburbs\textquotedblright{}.) Say it once more.
\end{tcolorbox}
